\section{Theory}
\label{sec:lab2_theory}

This section summarizes the signal theory needed for AM modulation and for the
two demodulation methods used in the lab. The emphasis is on the practical
signal-processing intuition behind the LabVIEW blocks rather than on formula
manipulation alone.

\subsection{Amplitude Modulation Fundamentals}

For standard AM with a transmitted carrier, the signal can be written as
\begin{equation}
s(t)=\left[A_c + A_m\,m(t)\right]\cos(2\pi f_c t),
\label{eq:lab2_am_general}
\end{equation}
where $m(t)$ is the baseband message, $A_c$ is the carrier amplitude, $A_m$ is
the message scaling amplitude, and $f_c$ is the carrier frequency.

For a single-tone message,
\begin{equation}
m(t)=\cos(2\pi f_m t),
\label{eq:lab2_single_tone}
\end{equation}
so the AM waveform becomes
\begin{equation}
s(t)=\left[A_c + A_m\cos(2\pi f_m t)\right]\cos(2\pi f_c t).
\label{eq:lab2_am_single_tone}
\end{equation}

The main parameters are:
\begin{itemize}
\item $A_c$: carrier amplitude,
\item $A_m$: message scaling amplitude,
\item $f_c$: carrier frequency,
\item $f_m$: message frequency in the single-tone case.
\end{itemize}

The modulation index is
\begin{equation}
\mu=\frac{A_m}{A_c}.
\label{eq:lab2_mu}
\end{equation}
The information is carried by the envelope
\begin{equation}
e(t)=A_c + A_m\,m(t)=A_c\left[1+\mu\,m(t)\right].
\end{equation}
In large-carrier AM, the carrier term is intentionally transmitted so that
simple envelope recovery remains possible.

\subsection{Spectral Characteristics of AM}

Let $M(f)$ denote the Fourier transform of $m(t)$. From
\eqref{eq:lab2_am_general},
\begin{equation}
S(f)=\frac{A_c}{2}\left[\delta(f-f_c)+\delta(f+f_c)\right]
+\frac{A_m}{2}\left[M(f-f_c)+M(f+f_c)\right].
\label{eq:lab2_am_spectrum_general}
\end{equation}
This immediately shows two effects:
\begin{itemize}
\item a carrier line at $\pm f_c$,
\item shifted copies of the baseband spectrum around $\pm f_c$.
\end{itemize}

For a single-tone message, the sidebands appear at
\begin{equation}
f_c-f_m \quad \text{and} \quad f_c+f_m.
\end{equation}
Hence the AM bandwidth is
\begin{equation}
B=2f_m
\end{equation}
for single-tone modulation and, more generally,
\begin{equation}
B=2B_m,
\end{equation}
where $B_m$ is the message bandwidth.

\subsection{Coherent Demodulation Theory}

Coherent demodulation multiplies the received signal by a synchronized local
carrier:
\begin{equation}
y(t)=s(t)\cos(2\pi f_c t+\phi),
\end{equation}
where $\phi$ is the receiver phase error.

Using $\cos\alpha\cos\beta=\tfrac{1}{2}[\cos(\alpha-\beta)+\cos(\alpha+\beta)]$,
\begin{equation}
y(t)=\frac{1}{2}\left[A_c+A_m m(t)\right]\cos\phi
+\frac{1}{2}\left[A_c+A_m m(t)\right]\cos(4\pi f_c t+\phi).
\label{eq:lab2_coherent_expand}
\end{equation}
Therefore:
\begin{itemize}
\item the first term is the desired baseband term,
\item the second term is centered around $2f_c$ and must be removed.
\end{itemize}

After low-pass filtering, the recovered signal is
\begin{equation}
y_{\mathrm{LPF}}(t)=\frac{1}{2}\left[A_c+A_m m(t)\right]\cos\phi.
\label{eq:lab2_coherent_lpf}
\end{equation}

Equation~\eqref{eq:lab2_coherent_lpf} directly explains the three main design
steps in the coherent detector:
\begin{itemize}
\item low-pass filtering removes the product term around $2f_c$,
\item DC removal is needed because the term proportional to $A_c$ appears as an
offset,
\item amplitude scaling is needed because the desired message is attenuated by
roughly a factor of $1/2$ in the ideal synchronized case.
\end{itemize}

\begin{tcolorbox}[breakable, colback=blue!5!white, colframe=blue!60!black,
title={\textbf{Intuition: Why Amplitude Scaling Is Needed}}]
After mixing the AM signal with the local oscillator (LO), the low-pass output
can be written in the more practical form
\begin{equation}
y_{\mathrm{LPF}}(t)=\frac{A_{\mathrm{LO}}}{2}\cos\phi
\left[A_c + A_m m(t)\right],
\end{equation}
where $A_{\mathrm{LO}}$ represents any LO amplitude gain in the receiver chain.
The attenuation therefore comes from two sources:
\begin{itemize}
\item the factor $\tfrac{1}{2}$ from the product-to-sum identity,
\item the factor $A_{\mathrm{LO}}\cos\phi$ from hardware gain and phase error.
\end{itemize}
In simulation, $A_{\mathrm{LO}}=1$ and $\phi=0$, so multiplying the recovered
message by $2$ restores the original amplitude. In hardware, the same idea
holds, but some calibration may still be needed because receiver gain is not
idealized.
\end{tcolorbox}

Coherence matters. If $\phi \neq 0$, the output is reduced by $\cos\phi$, and
if the local oscillator frequency is mismatched, the recovered waveform becomes
time-varyingly distorted instead of simply attenuated.

\subsection{Envelope Detection Theory}

Envelope detection is a simpler nonlinear alternative to coherent detection.
Conceptually, it consists of:
\begin{enumerate}
\item rectification,
\item low-pass filtering,
\item DC removal.
\end{enumerate}

For correct operation, the envelope must not invert. The
minimum envelope value is
\begin{equation}
e_{\min}=A_c-A_m=A_c(1-\mu).
\end{equation}
Therefore the no-inversion condition is
\begin{equation}
\mu \le 1.
\label{eq:lab2_no_inversion}
\end{equation}
When $\mu > 1$, over-modulation causes envelope inversion and the rectified,
filtered waveform no longer tracks the original message linearly.

Compared with coherent detection:
\begin{itemize}
\item coherent detection is linear when the carrier is properly synchronized and
therefore gives a more faithful recovery,
\item envelope detection is simpler to implement but is only approximate and is
more sensitive to amplitude impairments.
\end{itemize}

\subsection{Rectification Loss and Amplitude Scaling}

Rectification is nonlinear, so there is no single impulse response that applies
to all inputs. A useful approximation, and the one developed in the 
notes, is to model half-wave rectification as multiplication by a periodic
gating waveform $q(t)$ that keeps positive carrier half-cycles and suppresses
negative ones.

For a 50\% duty-cycle gate synchronized to the carrier, we may write
\begin{equation}
q(t)=\frac{1}{2}+q_{\mathrm{AC}}(t),
\end{equation}
where the AC part has zero mean and contains only odd harmonics of the carrier.
After multiplication and low-pass filtering, the recovered envelope is
attenuated by an approximately constant factor,
\begin{equation}
\hat{e}(t)\approx \frac{1}{\pi}\left[A_c + A_m m(t)\right].
\label{eq:lab2_pi_loss}
\end{equation}
This is why a factor of $\pi$ is often introduced in the envelope detector to
restore the expected amplitude before or after rectification, depending on the
implementation.

\begin{tcolorbox}[breakable, colback=orange!5!white,
colframe=orange!60!black,
title={\textbf{Intuition: Fourier Series of the Rectification Function (Part 1: Decomposition of $q(t)$)}}]
The half-wave rectifier switches between 1 and 0. With carrier period
$T_c = 1/f_c$ and angular carrier frequency $\omega_c = 2\pi f_c$, the
switching function is
\begin{equation}
q(t)=
\begin{cases}
1 & \cos(\omega_c t) \geq 0 \\
0 & \cos(\omega_c t) < 0.
\end{cases}
\end{equation}

We decompose this into a DC part and a zero-mean AC part:
\begin{equation}
q(t)=\frac{1}{2}+q_{\mathrm{AC}}(t),
\end{equation}

The DC term $\tfrac{1}{2}$ multiplied by $s(t)$ gives $\tfrac{1}{2}s(t)$, which
is still centered at $\pm f_c$ and is therefore removed completely by the
low-pass filter. So the DC part does not contribute to baseband recovery.

The AC part oscillates symmetrically between $+\tfrac{1}{2}$ and
$-\tfrac{1}{2}$:
\begin{equation}
q_{\mathrm{AC}}(t)=
\begin{cases}
+\tfrac{1}{2} & t \in [0,\,T_c/4) \cup [3T_c/4,\,T_c) \\
-\tfrac{1}{2} & t \in [T_c/4,\,3T_c/4).
\end{cases}
\end{equation}

Since $\cos(\omega_c t)$ is even, $q_{\mathrm{AC}}(t)$ is also even, so its
Fourier series contains only cosine terms and no DC:
\begin{equation}
q_{\mathrm{AC}}(t)=\sum_{n=1,3,5,\ldots} a_n \cos(n\omega_c t).
\end{equation}

Only odd harmonics appear because of half-wave symmetry:
\begin{equation}
q_{\mathrm{AC}}(t+T_c/2)=-q_{\mathrm{AC}}(t).
\end{equation}
\end{tcolorbox}

\begin{tcolorbox}[breakable, colback=orange!5!white,
colframe=orange!60!black,
title={\textbf{Intuition: Fourier Series of the Rectification Function (Part 2: Computing the Coefficients)}}]
\textbf{Step 1:} The cosine Fourier coefficients are
\begin{equation}
a_n = \frac{2}{T_c}\int_0^{T_c} q_{\mathrm{AC}}(t)\cos(n\omega_c t)\,dt.
\end{equation}

Splitting the integral over the intervals where
$q_{\mathrm{AC}}(t)=+\tfrac{1}{2}$ and $q_{\mathrm{AC}}(t)=-\tfrac{1}{2}$ gives
\begin{align}
a_n =\; &\frac{2}{T_c}\Bigg[
\frac{1}{2}\int_0^{T_c/4}\cos(n\omega_c t)\,dt \nonumber\\
&-\frac{1}{2}\int_{T_c/4}^{3T_c/4}\cos(n\omega_c t)\,dt \nonumber\\
&+\frac{1}{2}\int_{3T_c/4}^{T_c}\cos(n\omega_c t)\,dt\Bigg].
\end{align}

\textbf{Step 2:} Using
$\int \cos(n\omega_c t)\,dt = \frac{\sin(n\omega_c t)}{n\omega_c}$ and
$\omega_c T_c = 2\pi$,
\begin{align}
\int_0^{T_c/4}\cos(n\omega_c t)\,dt
&= \frac{\sin(n\pi/2)}{n\omega_c},\\[4pt]
\int_{T_c/4}^{3T_c/4}\cos(n\omega_c t)\,dt
&= \frac{\sin(3n\pi/2)-\sin(n\pi/2)}{n\omega_c},\\[4pt]
\int_{3T_c/4}^{T_c}\cos(n\omega_c t)\,dt
&= \frac{\sin(2n\pi)-\sin(3n\pi/2)}{n\omega_c}
= \frac{-\sin(3n\pi/2)}{n\omega_c}.
\end{align}

\textbf{Step 3:} Combining these terms and using $T_c\omega_c = 2\pi$ gives
\begin{align}
a_n &= \frac{1}{n\pi}\Bigg[
\sin\!\frac{n\pi}{2}
- \frac{1}{2}\!\left(\sin\!\frac{3n\pi}{2} - \sin\!\frac{n\pi}{2}\right)
- \frac{1}{2}\sin\!\frac{3n\pi}{2}
\Bigg] \nonumber\\[4pt]
&= \frac{1}{n\pi}\Bigg[
\frac{3}{2}\sin\!\frac{n\pi}{2}
- \sin\!\frac{3n\pi}{2}
\Bigg].
\end{align}

For odd $n$, $\sin(3n\pi/2)=-\sin(n\pi/2)$, so
\begin{equation}
\boxed{a_n = \frac{2\sin(n\pi/2)}{n\pi}}.
\end{equation}
\end{tcolorbox}

\begin{tcolorbox}[breakable, colback=orange!5!white,
colframe=orange!60!black,
title={\textbf{Intuition: Fourier Series of the Rectification Function (Part 3: Final Series and the $\tfrac{1}{\pi}$ Result)}}]
\textbf{Step 4:} Evaluating the coefficient formula,
\begin{itemize}
\item for even $n$, $\sin(n\pi/2)=0$, so $a_n=0$,
\item for odd $n$, the first coefficients are
\begin{align}
n=1:&\quad a_1 = +\frac{2}{\pi}, \\
n=3:&\quad a_3 = -\frac{2}{3\pi}, \\
n=5:&\quad a_5 = +\frac{2}{5\pi}.
\end{align}
\end{itemize}

Therefore
\begin{equation}
\boxed{q_{\mathrm{AC}}(t)=\frac{2}{\pi}\left[
\cos(\omega_c t)-\frac{1}{3}\cos(3\omega_c t)
+\frac{1}{5}\cos(5\omega_c t)-\cdots\right]}.
\end{equation}

The handout figure plots magnitudes of the spectral lines, which is why all
spikes appear positive even though the odd-harmonic coefficients alternate in
sign.

\textbf{Step 5:} The fundamental term is the one that produces the recovered
baseband. Multiplying
\begin{equation}
s(t)=\left[A_c + A_m m(t)\right]\cos(\omega_c t)
\end{equation}
by $\tfrac{2}{\pi}\cos(\omega_c t)$ gives
\begin{align}
[A_c + A_m m(t)]\cos(\omega_c t)\cdot \frac{2}{\pi}\cos(\omega_c t)
=\;&\frac{2}{\pi}[A_c + A_m m(t)]\cos^2(\omega_c t)\nonumber\\
=\;&\frac{2}{\pi}[A_c + A_m m(t)]\cdot\frac{1}{2}
\left[1+\cos(2\omega_c t)\right]\nonumber\\
=\;&\underbrace{\frac{1}{\pi}[A_c + A_m m(t)]}_{\text{baseband: kept by LPF}}
+\underbrace{\frac{1}{\pi}[A_c + A_m m(t)]\cos(2\omega_c t)}_{\text{at }2f_c\text{: removed by LPF}}.
\end{align}

After low-pass filtering,
\begin{equation}
\boxed{\hat{e}(t)=\frac{1}{\pi}[A_c + A_m m(t)]}.
\end{equation}

So the $\tfrac{1}{\pi}$ factor comes from the Fourier coefficient
$\tfrac{2}{\pi}$ of the rectification waveform and the extra $\tfrac{1}{2}$
from $\cos^2(\omega_c t)=\tfrac{1}{2}[1+\cos(2\omega_c t)]$. Multiplying by
$\pi$ in the LabVIEW VI compensates that dominant loss, while the remaining DC
term is later removed by the DC-removal stage.
\end{tcolorbox}

\begin{tcolorbox}[breakable, colback=purple!5!white,
colframe=purple!60!black,
title={\textbf{Intuition: Spectral Interpretation of Envelope Detection}}]
The same result can also be understood in the frequency domain. From the Fourier
series above, the rectification waveform produces a line spectrum. In magnitude
form, this can be written as
\begin{equation}
Q(f)=\frac{1}{2}\delta(f)
+\sum_{n=1,3,5,\ldots}\frac{1}{n\pi}
\left[\delta(f-nf_c)+\delta(f+nf_c)\right].
\end{equation}

Using the AM spectrum
\begin{equation}
S(f)=\frac{A_c}{2}\left[\delta(f-f_c)+\delta(f+f_c)\right]
+\frac{A_m}{2}\left[M(f-f_c)+M(f+f_c)\right],
\end{equation}
the rectified signal spectrum is
\begin{equation}
R(f)=S(f) * Q(f).
\end{equation}
Since $Q(f)$ is a sum of impulses at $\pm f_c$, $\pm 3f_c$, $\pm 5f_c$, and so
on, the convolution replicates the AM spectrum around each of these frequencies.
The impulses at $\pm f_c$ are the ones that shift copies of $S(f)$ down to
baseband, where they pass through the low-pass filter.
\end{tcolorbox}

\textbf{Why the factor is $\tfrac{2A}{\pi}$:} both the impulse at $+f_c$ and the
impulse at $-f_c$ have amplitude $\tfrac{A}{\pi}$, where $A=\tfrac{1}{2}$ is the
amplitude of the AC part of the rectification waveform. Each independently
shifts a copy of $S(f)$ to baseband, so their contributions add:
\begin{equation}
\frac{A}{\pi}+\frac{A}{\pi}=\frac{2A}{\pi}.
\end{equation}
Substituting $A=\tfrac{1}{2}$ gives
\begin{equation}
\frac{2A}{\pi}=\frac{1}{\pi},
\end{equation}
which matches the time-domain derivation above.

\begin{figure}[H]
\centering
\labincludegraphics[width=0.9\linewidth]{figures/convolutionrectifier.jpeg}
\caption{Rectification viewed as spectral replication of the AM signal.}
\label{fig:lab2_rectifier_convolution}
\end{figure}

\subsection{Role of Low-Pass Filtering}

Low-pass filtering is essential in both demodulators:
\begin{itemize}
\item in coherent detection it removes components near $2f_c$ and retains the
baseband term,
\item in envelope detection it smooths rectifier ripple and rejects higher
harmonics while retaining the slow envelope variation.
\end{itemize}

A practical cutoff rule is to choose the filter so that
\begin{equation}
f_{\mathrm{LP}} \gtrsim B_m
\quad \text{and} \quad
f_{\mathrm{LP}} \ll f_c.
\end{equation}
This allows the message to pass with low distortion while keeping most
carrier-related components out of the output.

Filter order also matters:
\begin{itemize}
\item higher order improves stopband rejection,
\item but higher order also increases phase delay and transient settling.
\end{itemize}
Butterworth filters are attractive in this lab because their passband is
maximally flat, which helps preserve message amplitude over the baseband range
of interest.

\subsection{Noise Effects in AM Systems}

Noise affects the two methods differently:
\begin{itemize}
\item envelope detection directly tracks amplitude, so amplitude noise appears
strongly at the output,
\item coherent detection is usually more robust because synchronized mixing and
low-pass filtering give better selectivity against out-of-band components.
\end{itemize}

In both cases, increasing $A_m$ and therefore increasing $\mu$ makes the
recovered message more visible. At very small $\mu$, the desired baseband can be
visually dominated by noise even when the AM signal is still present. For the
envelope detector, this improvement only remains meaningful while the
no-inversion condition $\mu \le 1$ is still satisfied.

\subsection{Approximate Noise Performance of the Two Demodulators}
\label{sec:lab2_noise_performance}

To connect the theory to the later experimental noise observations, let the
received AM signal be
\begin{equation}
r(t)=s(t)+n(t),
\end{equation}
where $n(t)$ is additive bandpass noise around the carrier. Around $f_c$, it is
convenient to write the noise in in-phase and quadrature form as
\begin{equation}
n(t)=n_I(t)\cos(2\pi f_c t)-n_Q(t)\sin(2\pi f_c t),
\end{equation}
with $n_I(t)$ and $n_Q(t)$ low-pass noise processes.

\subsubsection{Coherent Demodulation}

With perfect carrier synchronization, the receiver multiplies by
$\cos(2\pi f_c t)$ and then low-pass filters. The output becomes
\begin{equation}
y_c(t)\approx \frac{1}{2}\left[A_c+A_m m(t)\right]+\frac{1}{2}n_I(t).
\end{equation}
After DC removal, the useful part is
\begin{equation}
y_{c,\mathrm{sig}}(t)\approx \frac{A_m}{2}m(t),
\end{equation}
and the output noise is approximately
\begin{equation}
y_{c,\mathrm{noise}}(t)\approx \frac{1}{2}n_I(t).
\end{equation}

If the message has average power $P_m=\mathbb{E}[m^2(t)]$, then the coherent
detector output SNR is
\begin{equation}
\mathrm{SNR}_{c,\mathrm{out}}
\approx
\frac{(A_m^2/4)P_m}{\sigma_{n_I}^2/4}
=
\frac{A_m^2 P_m}{\sigma_{n_I}^2}.
\label{eq:lab2_snr_coherent}
\end{equation}
So, for fixed noise power and fixed carrier amplitude, the coherent-detector
output SNR grows proportionally to $A_m^2$, equivalently to $\mu^2$.

\subsubsection{Envelope Detection}

For envelope detection, assume the noise is small compared with the carrier and
the modulation stays within $\mu \le 1$. Then the noisy envelope can be
linearized as
\begin{equation}
\lvert r(t)\rvert
\approx
A_c+A_m m(t)+n_I(t).
\end{equation}
After low-pass filtering, DC removal, and amplitude correction, the baseband
output is approximately
\begin{equation}
y_e(t)\approx A_m m(t)+n_I(t).
\end{equation}
The corresponding first-order output SNR is therefore
\begin{equation}
\mathrm{SNR}_{e,\mathrm{out}}
\approx
\frac{A_m^2 P_m}{\sigma_{n_I}^2}.
\label{eq:lab2_snr_envelope}
\end{equation}

This shows that envelope detection has the same leading $\mu^2$ scaling under
high-SNR conditions. However, unlike coherent detection, it is nonlinear. Once
the carrier becomes weak relative to the noise, or when $\mu > 1$, neglected
higher-order terms become important and the envelope detector develops extra
distortion and a threshold-like breakdown. That is why coherent detection is
usually more robust in practice even though the first-order scaling laws look
similar.

\subsubsection{Interpretation for the Lab Results}

Equations~\eqref{eq:lab2_snr_coherent} and \eqref{eq:lab2_snr_envelope} imply
that for fixed noise level
\begin{equation}
\mathrm{SNR}_{\mathrm{out}} \propto \mu^2.
\end{equation}
Therefore:
\begin{itemize}
\item reducing $\mu$ from $1$ to $0.5$ reduces output SNR by a factor of $4$,
that is about $6$ dB,
\item reducing $\mu$ from $1$ to $0.2$ reduces output SNR by a factor of $25$,
that is about $14$ dB.
\end{itemize}
This matches the experimental observation that the received waveform becomes
progressively noisier as $\mu$ is reduced.

\subsection{Detailed Output SNR Models Under AWGN}
\label{sec:lab2_awgn_snr_models}

The earlier discussion established the main qualitative noise trends. The
following note now derives, step by step, the output SNR models for baseband
transmission, DSB-SC, and conventional AM under additive white Gaussian noise.
Throughout, the comparison reference is the \emph{baseband output SNR}
\begin{equation}
\left(\frac{S}{N}\right)_{\mathrm{baseband}}
=
\frac{P}{N_0 W},
\end{equation}
where $P=\mathbb{E}[m^2(t)]$ is the message power and the receiver output is
assumed bandlimited to the message band $\lvert f \rvert \leq W$.

\begin{tcolorbox}[breakable, colback=cyan!4!white, colframe=cyan!60!black,
title={\textbf{Intuition: Step-by-Step Output SNR Models Under AWGN}}]
\textbf{Common assumptions.}
\begin{itemize}
\item $m(t)$ is real, zero-mean, and bandlimited to $W$ Hz.
\item $P = \mathbb{E}[m^2(t)]$ denotes the message power.
\item $\omega_c = 2\pi f_c$ denotes carrier angular frequency.
\item The additive noise is AWGN with two-sided PSD $N_0/2$.
\item For the bandpass cases, after ideal bandpass selection around $\pm f_c$,
the narrowband noise is written as
\begin{equation}
n(t)=n_I(t)\cos(\omega_c t)-n_Q(t)\sin(\omega_c t),
\end{equation}
where $n_I(t)$ and $n_Q(t)$ are zero-mean low-pass Gaussian processes with
two-sided PSD $N_0/2$ over $\lvert f \rvert \leq W$.
\item Consequently, either low-pass noise component has power
\begin{equation}
\mathbb{E}[n_I^2(t)] = \mathbb{E}[n_Q^2(t)] = \int_{-W}^{W}\frac{N_0}{2}\,df = N_0 W.
\end{equation}
\end{itemize}

\textbf{1. Baseband transmission}

\emph{Transmitted signal model.}
\begin{equation}
s_{\mathrm{bb}}(t)=m(t).
\end{equation}
The transmitted signal power is therefore
\begin{equation}
\mathbb{E}[s_{\mathrm{bb}}^2(t)] = \mathbb{E}[m^2(t)] = P.
\end{equation}

\emph{Received signal model.}
\begin{equation}
r_{\mathrm{bb}}(t)=m(t)+n_{\mathrm{bb}}(t),
\end{equation}
where the baseband noise $n_{\mathrm{bb}}(t)$ has two-sided PSD $N_0/2$ over
$\lvert f \rvert \leq W$.

\emph{Demodulator model.}
For direct baseband transmission, the receiver is simply an ideal low-pass
output stage. Thus
\begin{equation}
y_{\mathrm{bb}}(t)=r_{\mathrm{bb}}(t)=m(t)+n_{\mathrm{bb}}(t).
\end{equation}

\emph{Output signal term.}
\begin{equation}
y_{\mathrm{bb, sig}}(t)=m(t).
\end{equation}
Its power is $P$ by definition.

\emph{Output noise term.}
\begin{equation}
y_{\mathrm{bb, noise}}(t)=n_{\mathrm{bb}}(t).
\end{equation}

\emph{Output noise power.}
\begin{equation}
P_{\mathrm{n,bb}}
=
\int_{-W}^{W}\frac{N_0}{2}\,df
=
N_0 W.
\end{equation}
This is the reference noise power over the message band.

\emph{Output SNR.}
\begin{equation}
\left(\frac{S}{N}\right)_{\mathrm{bb}}
=
\frac{P}{N_0 W}.
\label{eq:lab2_snr_baseband_reference}
\end{equation}
This is the baseline quantity against which all later schemes are compared.

\textbf{2. DSB-SC with coherent detection}

\emph{Transmitted signal model.}
\begin{equation}
s_{\mathrm{DSB}}(t)=m(t)\cos(\omega_c t).
\end{equation}

\emph{Received signal model.}
\begin{equation}
r_{\mathrm{DSB}}(t)=m(t)\cos(\omega_c t)+n(t).
\end{equation}
Using the narrowband noise representation,
\begin{equation}
r_{\mathrm{DSB}}(t)
=
m(t)\cos(\omega_c t)
+n_I(t)\cos(\omega_c t)-n_Q(t)\sin(\omega_c t).
\end{equation}

\emph{Demodulator model.}
To recover the message without amplitude loss, use a coherent local oscillator
$2\cos(\omega_c t)$ followed by an ideal low-pass filter:
\begin{equation}
y_{\mathrm{DSB}}(t)=\mathrm{LPF}\left\{2r_{\mathrm{DSB}}(t)\cos(\omega_c t)\right\}.
\end{equation}

\emph{Output signal term.}
First expand the signal part:
\begin{align}
2m(t)\cos(\omega_c t)\cos(\omega_c t)
&=
2m(t)\cos^2(\omega_c t)\nonumber\\
&=
m(t)\left[1+\cos(2\omega_c t)\right].
\end{align}
After low-pass filtering, the term at $2\omega_c$ is removed, so
\begin{equation}
\mathrm{LPF}\left\{2m(t)\cos^2(\omega_c t)\right\}=m(t).
\end{equation}
Thus the recovered signal term is exactly
\begin{equation}
y_{\mathrm{DSB, sig}}(t)=m(t).
\end{equation}

\emph{Output noise term.}
Now expand the noise contribution:
\begin{align}
2n(t)\cos(\omega_c t)
=\;&
2n_I(t)\cos^2(\omega_c t)
-2n_Q(t)\sin(\omega_c t)\cos(\omega_c t)\nonumber\\
=\;&
n_I(t)\left[1+\cos(2\omega_c t)\right]
-n_Q(t)\sin(2\omega_c t).
\end{align}
After low-pass filtering, only the low-frequency in-phase component remains:
\begin{equation}
y_{\mathrm{DSB, noise}}(t)=n_I(t).
\end{equation}

\emph{Output noise power.}
Since $n_I(t)$ has two-sided PSD $N_0/2$ over the message band,
\begin{equation}
P_{\mathrm{n,DSB}}
=
\mathbb{E}[n_I^2(t)]
=
N_0 W.
\end{equation}

\emph{Output SNR.}
\begin{equation}
\left(\frac{S}{N}\right)_{\mathrm{DSB-SC}}
=
\frac{P}{N_0 W}.
\end{equation}

\emph{Comparison with baseband.}
\begin{equation}
\left(\frac{S}{N}\right)_{\mathrm{DSB-SC}}
=
\left(\frac{S}{N}\right)_{\mathrm{baseband}}.
\end{equation}
So ideal coherent DSB-SC detection preserves the baseband SNR exactly.

\textbf{3. Conventional full AM with coherent detection}

\emph{Transmitted signal model.}
\begin{equation}
s_{\mathrm{AM}}(t)=\left[A_c+m(t)\right]\cos(\omega_c t).
\end{equation}

The total transmitted AM power is
\begin{align}
P_{\mathrm{T,AM}}
&=
\mathbb{E}[s_{\mathrm{AM}}^2(t)]\nonumber\\
&=
\mathbb{E}\!\left[(A_c+m(t))^2\cos^2(\omega_c t)\right]\nonumber\\
&=
\frac{1}{2}\mathbb{E}\!\left[A_c^2+2A_c m(t)+m^2(t)\right]\nonumber\\
&=
\frac{1}{2}\left(A_c^2+P\right),
\end{align}
because $\mathbb{E}[m(t)]=0$.

Therefore the carrier power is
\begin{equation}
P_c=\frac{A_c^2}{2},
\end{equation}
and the sideband power is
\begin{equation}
P_{\mathrm{SB}}=P_{\mathrm{T,AM}}-P_c=\frac{P}{2}.
\end{equation}
This already shows that only the sideband part is associated with the message.

\emph{Received signal model.}
\begin{equation}
r_{\mathrm{AM}}(t)=\left[A_c+m(t)\right]\cos(\omega_c t)+n(t).
\end{equation}

\emph{Demodulator model.}
Use the same coherent local oscillator and low-pass filter:
\begin{equation}
y_{\mathrm{AM,coh}}(t)=\mathrm{LPF}\left\{2r_{\mathrm{AM}}(t)\cos(\omega_c t)\right\}.
\end{equation}

\emph{Output signal term before DC removal.}
Expanding the desired part gives
\begin{align}
2\left[A_c+m(t)\right]\cos^2(\omega_c t)
&=
\left[A_c+m(t)\right]\left[1+\cos(2\omega_c t)\right].
\end{align}
After low-pass filtering,
\begin{equation}
\mathrm{LPF}\left\{2\left[A_c+m(t)\right]\cos^2(\omega_c t)\right\}
=
A_c+m(t).
\end{equation}
Thus the coherent detector output contains a DC term and the message term.

\emph{Output noise term.}
The bandpass noise goes through exactly the same algebra as in DSB-SC:
\begin{equation}
\mathrm{LPF}\left\{2n(t)\cos(\omega_c t)\right\}=n_I(t).
\end{equation}
So the output before DC removal is
\begin{equation}
y_{\mathrm{AM,coh}}(t)=A_c+m(t)+n_I(t).
\end{equation}

\emph{After DC removal.}
The carrier contribution $A_c$ carries no message information and is removed by
the DC-blocking stage. Therefore the useful recovered output is
\begin{equation}
\tilde{y}_{\mathrm{AM,coh}}(t)=m(t)+n_I(t).
\end{equation}

\emph{Output signal term.}
\begin{equation}
\tilde{y}_{\mathrm{AM,coh,sig}}(t)=m(t).
\end{equation}

\emph{Output noise term.}
\begin{equation}
\tilde{y}_{\mathrm{AM,coh,noise}}(t)=n_I(t).
\end{equation}

\emph{Output noise power.}
\begin{equation}
P_{\mathrm{n,AM,coh}}=N_0 W.
\end{equation}

\emph{Raw detector output SNR.}
\begin{equation}
\left(\frac{S}{N}\right)_{\mathrm{AM,coh,raw}}
=
\frac{P}{N_0 W}
=
\left(\frac{S}{N}\right)_{\mathrm{baseband}}.
\end{equation}
So, for a fixed message waveform $m(t)$, ideal coherent detection recovers the
same post-detection SNR as the baseband reference.

\emph{Power-efficiency form relative to total transmitted AM power.}
The subtle point is that conventional AM expends total transmitted power
\begin{equation}
P_{\mathrm{T,AM}}=\frac{A_c^2+P}{2},
\end{equation}
while only the sideband power
\begin{equation}
P_{\mathrm{SB}}=\frac{P}{2}
\end{equation}
is associated with the information-bearing component. Hence the useful fraction
of transmitted AM power is
\begin{equation}
\frac{P_{\mathrm{SB}}}{P_{\mathrm{T,AM}}}
=
\frac{P/2}{(A_c^2+P)/2}
=
\frac{P}{A_c^2+P}.
\end{equation}
Therefore, when the comparison is stated \emph{relative to total transmitted
AM power}, the standard coherent-AM relation is
\begin{equation}
\left(\frac{S}{N}\right)_{\mathrm{AM,coherent}}
=
\frac{P}{A_c^2+P}
\left(\frac{S}{N}\right)_{\mathrm{baseband}}.
\label{eq:lab2_awgn_am_coh_efficiency}
\end{equation}
The loss is entirely due to the carrier-power term $A_c^2/2$, which carries no
message information.

\textbf{4. Conventional full AM with envelope detection}

\emph{Transmitted signal model.}
\begin{equation}
s_{\mathrm{AM}}(t)=\left[A_c+m(t)\right]\cos(\omega_c t).
\end{equation}
The same AM power split derived above still applies:
\begin{equation}
P_{\mathrm{T,AM}}=\frac{A_c^2+P}{2},
\qquad
P_{\mathrm{SB}}=\frac{P}{2}.
\end{equation}

\emph{Received signal model.}
\begin{equation}
r_{\mathrm{AM}}(t)=\left[A_c+m(t)\right]\cos(\omega_c t)+n(t).
\end{equation}
Insert the narrowband noise decomposition:
\begin{align}
r_{\mathrm{AM}}(t)
=\;&
\left[A_c+m(t)\right]\cos(\omega_c t)
+n_I(t)\cos(\omega_c t)-n_Q(t)\sin(\omega_c t)\nonumber\\
=\;&
\left[A_c+m(t)+n_I(t)\right]\cos(\omega_c t)-n_Q(t)\sin(\omega_c t).
\end{align}

\emph{Envelope-detector model.}
The ideal envelope detector outputs the magnitude of the I/Q vector. Therefore,
up to this point, the following equation is exact:
\begin{equation}
v(t)=\sqrt{\left(A_c+m(t)+n_I(t)\right)^2+n_Q^2(t)}.
\label{eq:lab2_awgn_env_exact}
\end{equation}

\emph{Large-carrier / high-SNR approximation.}
Define
\begin{equation}
x(t)=A_c+m(t).
\end{equation}
Then \eqref{eq:lab2_awgn_env_exact} becomes
\begin{align}
v(t)
&=
\sqrt{\left(x(t)+n_I(t)\right)^2+n_Q^2(t)}\nonumber\\
&=
\sqrt{x^2(t)+2x(t)n_I(t)+n_I^2(t)+n_Q^2(t)}\nonumber\\
&=
x(t)\sqrt{1+\frac{2n_I(t)}{x(t)}
+\frac{n_I^2(t)+n_Q^2(t)}{x^2(t)}}.
\end{align}

Now let
\begin{equation}
u(t)=\frac{2n_I(t)}{x(t)}+\frac{n_I^2(t)+n_Q^2(t)}{x^2(t)}.
\end{equation}
If the carrier is strong, the post-detection SNR is high, and there is no
overmodulation so that $x(t)=A_c+m(t)$ stays positive and large, then
$\lvert u(t)\rvert \ll 1$. We may therefore use the first-order Taylor
expansion
\begin{equation}
\sqrt{1+u}\approx 1+\frac{u}{2}.
\end{equation}
Substituting this into the envelope expression gives
\begin{align}
v(t)
&\approx
x(t)\left[1+\frac{1}{2}u(t)\right]\nonumber\\
&=
x(t)\left[
1+\frac{n_I(t)}{x(t)}
+\frac{n_I^2(t)+n_Q^2(t)}{2x^2(t)}
\right].
\end{align}

Keeping only first-order terms in the noise means dropping the quadratic terms
$n_I^2(t)$ and $n_Q^2(t)$. Hence
\begin{equation}
v(t)\approx x(t)+n_I(t)=A_c+m(t)+n_I(t).
\label{eq:lab2_awgn_env_first_order}
\end{equation}
This is the key approximation step: it is not exact, and it is only valid for
high SNR, strong carrier, and no overmodulation.

\emph{After DC removal.}
Removing the large carrier-related DC term yields the approximate recovered
baseband model
\begin{equation}
\tilde{y}_{\mathrm{AM,env}}(t)\approx m(t)+n_I(t).
\end{equation}

\emph{Output signal term.}
\begin{equation}
\tilde{y}_{\mathrm{AM,env,sig}}(t)\approx m(t).
\end{equation}

\emph{Output noise term.}
\begin{equation}
\tilde{y}_{\mathrm{AM,env,noise}}(t)\approx n_I(t).
\end{equation}
To first order, only the in-phase noise component appears in the recovered
signal. The quadrature term $n_Q(t)$ enters only through the neglected
second-order terms.

\emph{Output noise power.}
\begin{equation}
P_{\mathrm{n,AM,env}}\approx N_0 W.
\end{equation}

\emph{Approximate output SNR.}
\begin{equation}
\left(\frac{S}{N}\right)_{\mathrm{AM,env,raw}}
\approx
\frac{P}{N_0 W}
=
\left(\frac{S}{N}\right)_{\mathrm{baseband}}.
\end{equation}
So, at first order and under strong-carrier conditions, the recovered
signal-plus-noise model is the same as the coherent result.

\emph{Power-efficiency form relative to total transmitted AM power.}
Because the transmitted conventional-AM power split is unchanged, the standard
power-efficiency statement is
\begin{equation}
\left(\frac{S}{N}\right)_{\mathrm{AM,env}}
\approx
\frac{P}{A_c^2+P}
\left(\frac{S}{N}\right)_{\mathrm{baseband}}.
\label{eq:lab2_awgn_am_env_efficiency}
\end{equation}

\emph{Figure of merit in terms of modulation index.}
For a single-tone message
\begin{equation}
m(t)=A_m\cos(\omega_m t),
\end{equation}
the message power is
\begin{equation}
P=\frac{A_m^2}{2}.
\end{equation}
If the modulation index is defined by
\begin{equation}
\mu=\frac{A_m}{A_c},
\end{equation}
then
\begin{equation}
P=\frac{A_c^2\mu^2}{2}.
\end{equation}
Substituting this into the efficiency factor gives
\begin{equation}
\frac{P}{A_c^2+P}
=
\frac{A_c^2\mu^2/2}{A_c^2+A_c^2\mu^2/2}
=
\frac{\mu^2}{2+\mu^2}.
\end{equation}
Hence the standard high-SNR envelope-detector figure of merit is
\begin{equation}
\left(\frac{S}{N}\right)_{\mathrm{AM,env}}
\approx
\frac{\mu^2}{2+\mu^2}
\left(\frac{S}{N}\right)_{\mathrm{baseband}}.
\end{equation}

\emph{Why envelope detection becomes worse at low carrier-to-noise ratio.}
The approximation in \eqref{eq:lab2_awgn_env_first_order} breaks down once the
carrier is no longer dominant. Then the neglected terms involving $n_Q^2(t)$,
$n_I^2(t)$, and the fluctuating denominator $A_c+m(t)$ are no longer small.
The envelope is no longer well approximated by a linear model, so the detector
develops distortion and a sharp threshold effect. This is why envelope detection
is usually inferior to coherent detection at low carrier-to-noise ratio, even
though their first-order high-SNR formulas look similar.
\end{tcolorbox}

\begin{table}[H]
\centering
\small
\renewcommand{\arraystretch}{1.2}
\begin{tabular}{|>{\raggedright\arraybackslash}p{0.16\linewidth}|>{\raggedright\arraybackslash}p{0.18\linewidth}|>{\raggedright\arraybackslash}p{0.28\linewidth}|>{\raggedright\arraybackslash}p{0.28\linewidth}|}
\hline
\textbf{Scheme} & \textbf{Demodulation} & \textbf{Output SNR} & \textbf{Relative efficiency / comments} \\
\hline
Baseband transmission & Direct baseband output & $\dfrac{P}{N_0 W}$ & Reference case. All transmitted power is useful message power. \\
\hline
DSB-SC & Ideal coherent detection & $\dfrac{P}{N_0 W}$ & Same output SNR as baseband. No carrier-power waste. \\
\hline
Conventional full AM & Coherent detection & Raw detector output: $\dfrac{P}{N_0 W}$. \newline Power-efficiency form: $\dfrac{P}{A_c^2+P}\left(\dfrac{S}{N}\right)_{\mathrm{baseband}}$ & Carrier power $A_c^2/2$ carries no information, so conventional AM is power-inefficient. \\
\hline
Conventional full AM & Envelope detection & First-order high-SNR result: $\dfrac{P}{N_0 W}$. \newline Power-efficiency form: $\approx \dfrac{P}{A_c^2+P}\left(\dfrac{S}{N}\right)_{\mathrm{baseband}}$ & Approximate result only. Valid for strong carrier, no overmodulation, and high CNR. Threshold effect causes clear degradation at low CNR. \\
\hline
\end{tabular}
\caption{Summary of output-SNR models under AWGN for the main AM-related schemes.}
\label{tab:lab2_awgn_snr_summary}
\end{table}
